\chapter{Fermion Masses: Chirality, Yukawa Couplings, and Majorana Neutrinos}

\lectureoverview{We first rebuild the gauge-boson Higgs mechanism in a toy
\(U(1)\) theory, then ask the parallel question for
fermions. A Dirac mass mixes left and right chiral fields, but the weak
interaction assigns them different charges and therefore forbids a bare
mass. A Yukawa coupling to the Higgs field respects the symmetry; after the
Higgs field acquires a vacuum value, the same interaction becomes a fermion
mass and a measurable Higgs--fermion coupling. The final part follows this
logic into neutrino mass, possible heavy unprotected particles, cosmic-ray
kinematics, relativistic inertia, and gravity.}

The purpose of the opening recap is not merely to repeat the previous
lecture. We need one idea from it in a form that can be carried over to
fermions: a mass is a nonzero energy for an excitation at zero momentum.
For a field, zero momentum means infinite wavelength, or a displacement that
is the same everywhere. A derivative-only theory cannot detect a uniform
displacement. A term without derivatives can.

\section{The Gauge-Boson Story in One Pass}

\subsection{Derivative energy and a massless vector}

The model gauge field is \(A_\mu\), with field strength
\begin{equation}
  F_{\mu\nu}
  =\partial_\mu A_\nu-\partial_\nu A_\mu .
  \label{eq:l09-field-strength}
\end{equation}
The board writes the gauge kinetic term schematically as \(F^2\). With a
standard normalization it is
\begin{equation}
  \mathcal L_{\mathrm{gauge}}
  =-\frac14 F_{\mu\nu}F^{\mu\nu}.
  \label{eq:l09-gauge-kinetic}
\end{equation}

\begin{figure}[htbp]
\centering
\includegraphics[width=0.72\textwidth]{lecture_09_figure_02.png}
\caption{The field-strength definition and the board's \(F^2\) shorthand
for the derivative energy of the gauge field.}
\lectureframe{00:02:39}
\label{fig:l09-field-strength-board}
\end{figure}

Every appearance of \(A_\mu\) in \eqref{eq:l09-gauge-kinetic} is
differentiated. Consequently,
\begin{equation}
  A_\mu(x)\longmapsto A_\mu(x)+c_\mu,
  \qquad \partial_\nu c_\mu=0,
\end{equation}
does not change \(F_{\mu\nu}\). A homogeneous vector potential therefore
costs no field-strength energy. In the language of the lecture, the
zero-momentum mode has no gap and the gauge boson is massless.

This \(U(1)\) field is a deliberately simple stand-in. The physical target
is the \(W^\pm\) and \(Z\), not the electromagnetic photon. The unbroken
electromagnetic symmetry leaves the photon massless; the broken
electroweak symmetry gives masses to the weak gauge bosons.

\subsection{Freezing the Higgs radius}

Introduce a charged complex scalar and write it in polar form,
\begin{equation}
  \phi(x)=\rho(x)e^{i\alpha(x)}.
\end{equation}
Its Mexican-hat potential is minimized on a circle of nonzero radius.
Calling that radius \(f\), the radial direction is comparatively stiff. For
the first pass we freeze it and retain only the phase:
\begin{equation}
  \phi(x)\simeq f e^{i\alpha(x)}.
  \label{eq:l09-frozen-higgs}
\end{equation}

\begin{figure}[htbp]
\centering
\includegraphics[width=0.72\textwidth]{lecture_09_figure_03.png}
\caption{The frozen Higgs radius, its covariant derivative, and the small
broken-potential sketch used in the recap.}
\lectureframe{00:06:42}
\label{fig:l09-frozen-board}
\end{figure}

Choose the convention
\begin{equation}
\begin{aligned}
  D_\mu&=\partial_\mu+igA_\mu,
  \\
  \phi&\longmapsto e^{i\theta}\phi,
  \\
  A_\mu&\longmapsto A_\mu-\frac1g\partial_\mu\theta .
\end{aligned}
  \label{eq:l09-gauge-convention}
\end{equation}
The lecture suppresses \(g\); keeping it visible makes the dimensions and
the eventual mass transparent. Acting on the frozen field gives
\begin{align}
  D_\mu\phi
  &=(\partial_\mu+igA_\mu)f e^{i\alpha}
  \notag\\
  &=i\bigl(\partial_\mu\alpha+gA_\mu\bigr)
    f e^{i\alpha}.
  \label{eq:l09-covariant-frozen}
\end{align}
Multiplication by the complex conjugate removes the phase and also turns
\(i\) times \(-i\) into \(+1\):
\begin{equation}
  |D_\mu\phi|^2
  =f^2
   \bigl(\partial_\mu\alpha+gA_\mu\bigr)
   \bigl(\partial^\mu\alpha+gA^\mu\bigr).
  \label{eq:l09-frozen-square}
\end{equation}

\begin{classroomqa}
\audiencequestion{Should the square of the covariant derivative carry a
minus sign because the derivative produced a factor of \(i\)?
\lecturetimestamp{00:07:39}}
\lecturerresponse{No. The kinetic term is the product with the complex
conjugate. One factor contributes \(i\), the other \(-i\), and
\(i(-i)=1\). The sign convention for \(D_\mu\) may be reversed consistently,
but the modulus square is positive in the corresponding energy expression.}
\end{classroomqa}

The combination in \eqref{eq:l09-frozen-square} is itself a shifted gauge
field. Choose \(\theta=-\alpha\), so that the scalar phase disappears, and
define
\begin{equation}
  A'_\mu=A_\mu+\frac1g\partial_\mu\alpha .
  \label{eq:l09-unitary-gauge-field}
\end{equation}
Then
\begin{equation}
  |D_\mu\phi|^2=g^2f^2 A'_\mu A'^\mu .
  \label{eq:l09-vector-mass-term}
\end{equation}
The new term contains no derivative of \(A'_\mu\). A homogeneous field now
stores energy, so the zero-momentum vector mode is gapped.

\begin{figure}[htbp]
\centering
\includegraphics[width=0.72\textwidth]{lecture_09_figure_04.png}
\caption{The board isolates the \(f^2A^2\) term that signals a gauge-boson
mass after the phase has been absorbed.}
\lectureframe{00:10:31}
\label{fig:l09-vector-mass-board}
\end{figure}

With canonical scalar normalization one writes
\(\phi=(f+h)e^{i\pi/f}/\sqrt2\). Comparing the resulting vector term with
\(\tfrac12m_A^2A_\mu A^\mu\) gives
\begin{equation}
  m_A=gf
  \label{eq:l09-abelian-vector-mass}
\end{equation}
in this abelian model. The lecture suppresses the normalization factors and
summarizes the result by saying that the vector mass is set by \(f\).

The phase mode has not vanished without accounting. Before gauging, it was
a massless Goldstone boson. After gauging, it supplies the longitudinal
polarization needed by a massive vector particle. The radial fluctuation,
which we temporarily froze, remains as the physical Higgs excitation.

\subsection{Why a term without derivatives means mass}

\begin{classroomqa}
\audiencequestion{Why does an \(A'^2\) term count as a mass rather than just
another interaction?\lecturetimestamp{00:13:10}}
\lecturerresponse{Mass is energy at rest. For a field, a particle at rest is
the quantum of a zero-momentum, infinite-wavelength mode. The derivative
term gives no energy to a static homogeneous displacement, whereas
\(m_A^2A'^2\) does. That nonzero zero-momentum energy is the rest-energy gap.}
\end{classroomqa}

The lecture develops this statement through three mechanical analogies.

\begin{itemize}
  \item In a crystal, slowly varying displacement costs elastic energy
  because neighboring atoms are shifted differently. A rigid translation
  shifts every atom together and costs no internal energy. The acoustic
  phonon is therefore gapless.
  \item In an ideal rotationally symmetric magnet, a spatially varying spin
  direction costs exchange energy, but rotating every spin together does
  not. The associated long-wavelength spin wave is gapless.
  \item In a plasma, shifting the electron fluid relative to the positive
  background creates charge separation and a restoring electric field even
  for a homogeneous relative displacement. The resulting plasmon has a
  nonzero frequency at zero wave number.
\end{itemize}

The Higgsed gauge field behaves like the third example: the broken medium
resists what would otherwise be a cost-free homogeneous displacement.

\begin{classroomqa}
\audiencequestion{Before the gauge field was introduced, did the phase
\(\alpha\) not already vibrate?\lecturetimestamp{00:18:24}}
\lecturerresponse{It did. Its energy was proportional to
\(f^2(\partial_\mu\alpha)^2\). A phase that varies from place to place has
gradient energy and can propagate, but a uniform change of phase costs
nothing. That is precisely the massless Goldstone mode.}
\end{classroomqa}

\begin{classroomqa}
\audiencequestion{The derivative includes a time derivative. Does actually
moving the whole lattice or changing the whole field not carry kinetic
energy?\lecturetimestamp{00:22:18}}
\lecturerresponse{Yes. The comparison is between static configurations, or
between configurations prepared arbitrarily slowly. A lattice translating
with nonzero velocity has kinetic energy. The mass test asks whether the
configuration still has energy when it is at rest and spatially uniform.}
\end{classroomqa}

\section{The Experimental Clue: Weak Interactions Distinguish Left and Right}

The lecture now turns to the electron, muon, and quarks. Their masses also
come from the Higgs field, but the mechanism looks different because a
fermion mass is not a quadratic vector term. It is a coupling between two
chiral fields.

\subsection{Parity violation in beta decay}

Begin with the empirical fact that the weak interaction violates parity.
Quantum electrodynamics and quantum chromodynamics treat a process and its
mirror image alike. The weak interaction does not. In neutron beta decay,
\begin{equation}
  n\longrightarrow p+e^-+\bar\nu_e,
  \label{eq:l09-beta-decay}
\end{equation}
the high-energy electrons selected by the charged-current interaction have
left-handed helicity: their spin points opposite to their momentum. The
corresponding antineutrinos have the opposite helicity.

For an ultrarelativistic particle, define helicity by
\begin{equation}
  h=\frac{\bm S\cdot\bm p}{|\bm S|\,|\bm p|}.
\end{equation}
Spin parallel to momentum has \(h=+1\); spin antiparallel has \(h=-1\).
The weak interaction couples to a left-chiral electron field and a
left-chiral neutrino field, not to a right-chiral electron field in the same
way.\editorialnote{The lecture initially uses handedness in the helicity sense
because that is visually intuitive at high energy. Gauge charges actually
act on chirality, the Lorentz-representation label selected by
\(P_L=(1-\gamma^5)/2\) and \(P_R=(1+\gamma^5)/2\). For a massless particle,
chirality and helicity coincide for particles; for a massive particle they
do not.}

To expose the obstruction without the full
\(SU(2)_L\times U(1)_Y\) bookkeeping, imagine a fake \(U(1)\) world in which
only \(\psi_L\) carries the relevant charge:
\begin{equation}
  \psi_L\longmapsto e^{i\theta}\psi_L,
  \qquad
  \psi_R\longmapsto\psi_R .
  \label{eq:l09-chiral-toy-charge}
\end{equation}
This is not electromagnetism: real left- and right-chiral electrons carry
the same electric charge. It is a caricature of the weak interaction, in
which the two chiralities sit in different electroweak representations.

\begin{classroomqa}
\audiencequestion{Does the same left--right asymmetry apply to quarks?
\lecturetimestamp{00:31:50}}
\lecturerresponse{Yes, with the appropriate weak multiplets and
hypercharges. Left-chiral up- and down-type quarks form weak doublets,
whereas the right-chiral fields are weak singlets. Their Higgs couplings
therefore face the same structural problem as the electron's.}
\end{classroomqa}

\section{A Dirac Mass Is Left--Right Mixing}

\subsection{From the four-component equation to two Weyl equations}

A Dirac wavefunction has four components. Loosely, two describe spin and
two distinguish the positive- and negative-energy solutions, with the latter
reinterpreted as antiparticle degrees of freedom. In Hamiltonian form,
\begin{equation}
  i\frac{\partial\Psi}{\partial t}
  =\left(-i\bm\alpha\cdot\bm\nabla+m\beta\right)\Psi .
  \label{eq:l09-dirac-hamiltonian}
\end{equation}
The matrices \(\alpha_i\) and \(\beta\) encode the relativistic relation
between energy, momentum, and mass.

\begin{figure}[htbp]
\centering
\includegraphics[width=0.74\textwidth]{lecture_09_figure_05.png}
\caption{The board first writes the Dirac equation, then splits it into
coupled left- and right-handed equations.}
\lectureframe{00:40:30}
\label{fig:l09-dirac-board}
\end{figure}

The chiral basis makes the role of \(m\) explicit. With
\(\sigma^\mu=(1,\bm\sigma)\) and
\(\bar\sigma^\mu=(1,-\bm\sigma)\),
\begin{align}
  i\bar\sigma^\mu\partial_\mu\psi_L&=m\psi_R,
  \label{eq:l09-weyl-left}\\
  i\sigma^\mu\partial_\mu\psi_R&=m\psi_L.
  \label{eq:l09-weyl-right}
\end{align}
At \(m=0\), the equations are independent: a left-chiral field and a
right-chiral field can propagate without talking to each other. At
\(m\neq0\), each appears as the source for the other. This is the exact
sense in which a Dirac mass mixes left and right.

Susskind describes this as the fermion oscillating between left and right.
That language captures the coupling but should not be pushed into a literal
classical picture of a point particle flipping back and forth. The precise
statement is the coupled field equation
\eqref{eq:l09-weyl-left}--\eqref{eq:l09-weyl-right}.

\begin{classroomqa}
\audiencequestion{Are the factors of \(i\) and the signs in the two board
equations consistent?\lecturetimestamp{00:41:10}}
\lecturerresponse{Not line by line on the board; an \(i\) was pulled outside
inconsistently during the quick derivation. The reliable content is the
coupling pattern. Equations \eqref{eq:l09-weyl-left} and
\eqref{eq:l09-weyl-right} give one standard consistent convention.}
\end{classroomqa}

\subsection{Why the weak symmetry forbids a bare mass}

The corresponding Dirac mass term in the Lagrangian is
\begin{equation}
  \mathcal L_m
  =-m\bar\psi\psi
  =-m\left(\bar\psi_L\psi_R+\bar\psi_R\psi_L\right).
  \label{eq:l09-bare-dirac-mass}
\end{equation}
Under the toy transformation \eqref{eq:l09-chiral-toy-charge},
\begin{equation}
  \bar\psi_L\psi_R\longmapsto
  e^{-i\theta}\bar\psi_L\psi_R,
  \qquad
  \bar\psi_R\psi_L\longmapsto
  e^{i\theta}\bar\psi_R\psi_L.
\end{equation}
Neither term is invariant. The direct left--right mixing coefficient \(m\)
is therefore forbidden before symmetry breaking.

Electromagnetism alone would not cause this problem, because \(\psi_L\) and
\(\psi_R\) have the same electric charge. The obstruction arises from their
different weak quantum numbers.

\subsection{Helicity questions that expose the chirality subtlety}

The class now tests the intuitive handedness picture from several angles.
These questions are not side issues: they reveal exactly where helicity is
useful and where chirality must replace it.

\begin{classroomqa}
\audiencequestion{What would the world be like if the weak interaction
emitted both electron helicities equally?\lecturetimestamp{00:42:12}}
\lecturerresponse{The immediate parity asymmetry would disappear from that
process. But the real electroweak theory is chiral: only the left-chiral
doublets carry the charged weak interaction. The lecture postpones a fuller
account of why nature chose that structure.}
\end{classroomqa}

\begin{classroomqa}
\audiencequestion{Suppose an electron's spin is prepared vertically and the
electron is then accelerated horizontally. Is it left- or right-handed?
\lecturetimestamp{00:42:44}}
\lecturerresponse{A spin state transverse to the momentum is a superposition
of the two helicities. Helicity is sharp only when the spin is parallel or
antiparallel to the momentum. The weak chiral projectors are nevertheless
well defined for the field.}
\end{classroomqa}

A massive particle can be slowed to rest and accelerated in the opposite
direction without reversing its spin. Its helicity then changes sign. A
massless particle cannot be overtaken or brought to rest by a Lorentz boost,
so its helicity is invariant under proper boosts. This motivates the
lecture's statement that a species with only one physical helicity must be
massless. The more exact field-theory statement is that a Lorentz-invariant
Dirac mass requires both chiralities.

\begin{classroomqa}
\audiencequestion{A circularly polarized photon can reflect from a mirror
and reverse direction without ever coming to rest. Does that contradict the
helicity argument?\lecturetimestamp{00:45:28}}
\lecturerresponse{No. The mirror is an external system that can exchange
angular momentum and reverse the propagation direction. A reflection
process is not a Lorentz boost that overtakes the photon. Whether a channel
exists also depends on the allowed polarization states and on the material
boundary conditions.}
\end{classroomqa}

\begin{classroomqa}
\audiencequestion{Does an ordinary electron beam contain one handedness or
both?\lecturetimestamp{00:46:10}}
\lecturerresponse{Both in general. A massive electron state contains both
chiral components, and an unpolarized beam contains both helicities.
Specialized sources can produce a polarized beam with a large helicity
asymmetry.}
\end{classroomqa}

\begin{classroomqa}
\audiencequestion{Can a beam be focused independently of its spin, or do
magnetic fields separate the two polarizations?\lecturetimestamp{00:46:29}}
\lecturerresponse{Ordinary position focusing can be largely insensitive to
spin, but magnetic moments couple to magnetic-field gradients. Spin-selective
sources, Stern--Gerlach-like gradients, and accelerator spin dynamics can
therefore prepare or separate polarizations.}
\end{classroomqa}

\begin{classroomqa}
\audiencequestion{Is the weak difference between left and right merely an
accident?\lecturetimestamp{00:48:13}}
\lecturerresponse{For the present course it is an experimental fact built
into the Standard Model. It is too systematic to feel like a casual
accident, but the Standard Model itself does not derive the complete pattern
of chiral charges from a deeper principle.}
\end{classroomqa}

\section{The Yukawa Construction}

\subsection{Let the Higgs carry the missing charge}

The symmetry forbids a numerical mass, but it allows an interaction with a
field that supplies the missing transformation phase. Introduce a complex
scalar with
\begin{equation}
  \phi\longmapsto e^{i\theta}\phi,
  \qquad
  \phi^\dagger\longmapsto e^{-i\theta}\phi^\dagger.
  \label{eq:l09-scalar-charge}
\end{equation}
Then
\begin{equation}
  \mathcal L_Y
  =-y\,\bar\psi_L\phi\psi_R
   -y^*\,\bar\psi_R\phi^\dagger\psi_L
  \label{eq:l09-yukawa-lagrangian}
\end{equation}
is gauge invariant. In the first term,
\(\bar\psi_L\) contributes \(e^{-i\theta}\) and \(\phi\) contributes
\(e^{i\theta}\); the phases cancel. The second term is its Hermitian
conjugate.

Equivalently, the chiral equations acquire field-dependent mixing:
\begin{align}
  i\bar\sigma^\mu D_\mu\psi_L
    &=y\phi\psi_R,
  \label{eq:l09-yukawa-eom-left}\\
  i\sigma^\mu\partial_\mu\psi_R
    &=y^*\phi^\dagger\psi_L.
  \label{eq:l09-yukawa-eom-right}
\end{align}

\begin{figure}[htbp]
\centering
\includegraphics[width=0.76\textwidth]{lecture_09_figure_06.png}
\caption{The bare mass on the right-hand sides has been replaced by the
charged scalar field. The exact board signs are schematic; the invariant
structure is given in \eqref{eq:l09-yukawa-lagrangian}.}
\lectureframe{00:51:30}
\label{fig:l09-yukawa-board}
\end{figure}

Before symmetry breaking, this is an interaction rather than a mass. A
right-chiral fermion can become a left-chiral fermion only while emitting or
absorbing the scalar quantum required to balance the charge. The coefficient
\(y\) is a Yukawa coupling, not the gauge charge \(g\).

\begin{classroomqa}
\audiencequestion{Is it legitimate simply to put \(\phi\) and
\(\phi^\dagger\) on the right-hand sides of the two equations?
\lecturetimestamp{00:51:55}}
\lecturerresponse{It is legitimate precisely because their transformation
phases repair the mismatch. The decisive test is not the visual form of the
equation but gauge invariance of the interaction
\(\bar\psi_L\phi\psi_R+\mathrm{h.c.}\).}
\end{classroomqa}

\subsection{The vacuum converts interaction into mass}

Now identify \(\phi\) with the Higgs field and let it settle at nonzero
vacuum magnitude. In the lecture's simplified normalization,
\begin{equation}
  \phi\simeq f,
\end{equation}
so \eqref{eq:l09-yukawa-lagrangian} contains
\begin{equation}
  \mathcal L_Y
  \supset
  -yf\,\bar\psi_L\psi_R
  -y^*f\,\bar\psi_R\psi_L .
\end{equation}
After a harmless phase redefinition when appropriate, the fermion mass is
\begin{equation}
  m_f=yf.
  \label{eq:l09-lecture-fermion-mass}
\end{equation}
This is the fermionic counterpart of the vector result \(m_A=gf\). The same
Higgs vacuum that supplies a longitudinal polarization to the weak gauge
bosons supplies the missing weak quantum numbers needed to join left and
right fermions.

In the canonically normalized Standard Model convention,
\begin{equation}
  H(x)=\frac1{\sqrt2}
  \begin{pmatrix}
    0\\ v+h(x)
  \end{pmatrix},
  \qquad v\simeq246\ \mathrm{GeV},
  \label{eq:l09-sm-higgs-vacuum}
\end{equation}
and therefore
\begin{equation}
  m_f=\frac{y_fv}{\sqrt2}.
  \label{eq:l09-sm-fermion-mass}
\end{equation}
The lecture's \(f\) and the conventional \(v/\sqrt2\) differ only by this
choice of normalization.

\begin{classroomqa}
\audiencequestion{If only the product \(yf\) is a fermion mass, how can the
coupling \(y\) and vacuum scale \(f\) be determined separately?
\lecturetimestamp{00:55:36}}
\lecturerresponse{The weak gauge-boson masses determine the vacuum scale once
the gauge couplings are known. A fermion's measured mass then determines its
Yukawa coupling. In the full theory,
\(m_W=gv/2\), \(m_Z=v\sqrt{g^2+g'^2}/2\), and
\(y_f=\sqrt2m_f/v\).}
\end{classroomqa}

\begin{classroomqa}
\audiencequestion{Is the weak-boson mass about \(900\ \mathrm{GeV}\)?
\lecturetimestamp{00:57:43}}
\lecturerresponse{No; the lecture immediately corrects the slip to roughly
\(90\ \mathrm{GeV}\). More precisely,
\(m_W\simeq80\ \mathrm{GeV}\) and \(m_Z\simeq91\ \mathrm{GeV}\), while the
underlying vacuum scale is \(v\simeq246\ \mathrm{GeV}\).}
\end{classroomqa}

\subsection{One vacuum scale, many Yukawa couplings}

Every charged fermion species has its own Yukawa coupling:
\begin{equation}
  m_i=y_i f
  \qquad\text{(lecture normalization)}.
  \label{eq:l09-mass-family}
\end{equation}

\begin{figure}[htbp]
\centering
\includegraphics[width=0.76\textwidth]{lecture_09_figure_08.png}
\caption{The board summarizes the species dependence as \(m_i=g_i f\);
the \(g_i\) written there are the Yukawa couplings called \(y_i\) here.}
\lectureframe{01:16:30}
\label{fig:l09-mass-family-board}
\end{figure}

For the electron, the lecture estimates
\begin{equation}
  y_e\sim\frac{0.5\ \mathrm{MeV}}{10^2\ \mathrm{GeV}}
  \sim10^{-5}.
\end{equation}
Using the canonical value \(v=246\ \mathrm{GeV}\) gives the more precise
\(y_e\simeq2.9\times10^{-6}\). The top quark lies at the opposite end:
\begin{equation}
  y_t=\frac{\sqrt2m_t}{v}\simeq1.
\end{equation}
Thus the Higgs mechanism explains how a gauge-invariant mass can exist, but
it does not explain the enormous hierarchy among the \(y_i\). The observed
fermion masses merely translate that hierarchy into measured numbers.

\begin{classroomqa}
\audiencequestion{Does the top quark really have a Yukawa coupling of order
one?\lecturetimestamp{00:59:09}}
\lecturerresponse{Yes. Its mass is comparable to the electroweak scale, so
its coupling is not a tiny perturbative number like the electron's. In the
modern normalization \(y_t\) is close to unity.}
\end{classroomqa}

\subsection{Restoring the physical Higgs fluctuation}

Freezing the scalar at its vacuum value hid the radial excitation. Restore
it schematically:
\begin{equation}
  \phi(x)\simeq f+h(x).
\end{equation}
The Yukawa factor becomes
\begin{equation}
  y\phi\simeq yf+yh=m_f+yh.
  \label{eq:l09-mass-plus-higgs}
\end{equation}

\begin{figure}[htbp]
\centering
\includegraphics[width=0.76\textwidth]{lecture_09_figure_07.png}
\caption{The board restores the Higgs fluctuation inside the Yukawa term:
the constant part produces the fermion mass and the fluctuating part
produces Higgs emission or absorption.}
\lectureframe{01:01:15}
\label{fig:l09-higgs-fluctuation-board}
\end{figure}

The first term in \eqref{eq:l09-mass-plus-higgs} is the mass. The second is a
vertex in which a Higgs boson couples to the fermion. In canonical Standard
Model notation,
\begin{equation}
  \mathcal L_Y
  \supset
  -m_f\bar\psi\psi
  -\frac{m_f}{v}h\bar\psi\psi.
  \label{eq:l09-higgs-fermion-vertex}
\end{equation}
The coupling is proportional to the fermion mass. Subject to kinematics,
color factors, phase space, and competing channels, the Higgs therefore
couples more strongly to heavier fermions.

\begin{classroomqa}
\audiencequestion{Are coupling constants themselves probabilities?
\lecturetimestamp{01:03:35}}
\lecturerresponse{They enter quantum amplitudes. A simple rate is often
proportional to the modulus square of an amplitude and hence to a coupling
squared, but phase space, spin sums, interference, loops, and multiple
vertices also matter. A coupling should not be identified directly with a
probability.}
\end{classroomqa}

\begin{classroomqa}
\audiencequestion{Does this argument assume that there is only one Higgs
field?\lecturetimestamp{01:03:59}}
\lecturerresponse{It describes the minimal Standard Model, which has one
complex Higgs doublet. Three of its four real components are absorbed by
\(W^\pm\) and \(Z\), leaving one physical neutral scalar. Extensions can
contain more doublets or other scalar multiplets. Supersymmetric models
typically require two Higgs doublets and therefore contain several physical
Higgs states.}
\end{classroomqa}

That historical expectation can now be compared with experiment.
\editorialnote{This lecture was recorded before Higgs discovery. ATLAS and
CMS announced a Higgs-like particle in 2012, now measured near
\(125\ \mathrm{GeV}\). Its observed couplings are broadly consistent with
the mass-proportional Standard Model pattern. The statement that it decays
preferentially to the heaviest particles always carries the kinematic
caveat: an on-shell \(t\bar t\) pair is too heavy for a \(125\ \mathrm{GeV}\)
Higgs.}

\section{Antiparticles, Beams, and the Meaning of Chirality}

Before turning to neutrinos, the class returns to several operational
questions about antiparticles and accelerator beams.

\begin{classroomqa}
\audiencequestion{What handedness does the antiparticle of a left-handed
electron have?\lecturetimestamp{01:06:05}}
\lecturerresponse{A left-chiral electron field annihilates left-helicity
electrons and creates right-helicity positrons in the ultrarelativistic
limit. This is one place where particle versus antiparticle and chirality
versus helicity must be stated carefully.}
\end{classroomqa}

\begin{classroomqa}
\audiencequestion{Can an electron and positron collide to make a \(Z\)
boson?\lecturetimestamp{01:07:13}}
\lecturerresponse{Yes, when their center-of-mass energy matches the allowed
channel. Electron--positron colliders studied the \(Z\) resonance with great
precision. Charge and the other quantum numbers are conserved by the full
process.}
\end{classroomqa}

\begin{classroomqa}
\audiencequestion{Are positrons all produced through one special weak
process?\lecturetimestamp{01:07:37}}
\lecturerresponse{No. Pair production is a standard electromagnetic route.
For example, an energetic electron can radiate a virtual photon, and that
photon can produce an \(e^-e^+\) pair in the field of another particle. The
lecture sketches this with a Feynman diagram.}
\end{classroomqa}

\begin{classroomqa}
\audiencequestion{Can a linear accelerator select the speed or energy of its
electrons?\lecturetimestamp{01:09:24}}
\lecturerresponse{Yes. The accelerating fields and timing select an energy
range. Once electrons are ultrarelativistic their speed is already extremely
close to \(c\), so additional accelerator energy mostly increases momentum
and total energy rather than visibly increasing speed.}
\end{classroomqa}

\begin{classroomqa}
\audiencequestion{Can an accelerator prepare a beam of one polarization?
\lecturetimestamp{01:09:47}}
\lecturerresponse{Yes, although it is technically demanding. One can prepare
spin-polarized electrons at low energy with a spin-dependent source and then
accelerate them while controlling spin precession. Modern colliders quantify
the resulting beam polarization rather than treating it as perfectly pure.}
\end{classroomqa}

\section{Neutrino Mass and the Majorana Possibility}

\subsection{Why neutrino mass forced an extension}

The original minimal Standard Model included only left-chiral neutrino
fields and therefore predicted massless neutrinos. Neutrino oscillations
show that at least two neutrino mass eigenvalues differ and that neutrino
masses are nonzero, although tiny compared with charged-fermion masses.

\begin{classroomqa}
\audiencequestion{How small are the neutrino masses, and where is their
right-handed component?\lecturetimestamp{01:10:38}}
\lecturerresponse{Oscillations determine mass-squared differences, not the
absolute masses. The relevant splittings correspond to scales of roughly
\(0.009\ \mathrm{eV}\) and \(0.05\ \mathrm{eV}\). A Dirac description adds
sterile right-chiral neutrinos; a Majorana description instead relates the
missing chirality to charge-conjugate neutrino fields.}
\end{classroomqa}

For a Dirac neutrino one adds a new field \(\nu_R\) and repeats the Yukawa
construction:
\begin{equation}
  \mathcal L_{\nu,\mathrm{Dirac}}
  =-y_\nu\bar L\,\widetilde H\,\nu_R+\mathrm{h.c.},
  \qquad
  m_\nu=\frac{y_\nu v}{\sqrt2}.
  \label{eq:l09-dirac-neutrino}
\end{equation}
The required \(y_\nu\) is extraordinarily small.

There is another logical possibility because the neutrino carries no
electric charge. A left-chiral neutrino can be coupled to the charge
conjugate of itself, which has the opposite chirality:
\begin{equation}
  \nu_L\longleftrightarrow(\nu_L)^c_R .
  \label{eq:l09-majorana-mixing}
\end{equation}
A mass eigenstate built this way is a Majorana fermion: it is its own
antiparticle, up to phase conventions. The analogous electron--positron
mixing is forbidden because it would violate exact electric-charge
conservation.

\begin{classroomqa}
\audiencequestion{Why is particle--antiparticle mixing allowed for a
neutrino but not for an electron?\lecturetimestamp{01:13:43}}
\lecturerresponse{An electron and a positron have opposite electric charge,
so a mass term that mixes them as the same one-particle state would violate
the unbroken electromagnetic symmetry. A neutrino is electrically neutral.
If no exact conserved lepton number distinguishes it from its antiparticle,
a Majorana mass is allowed.}
\end{classroomqa}

The name honors Ettore Majorana, the young Italian physicist who formulated
this kind of real fermion representation and disappeared in 1938. The
lecture pauses over the historical mystery: his fate was never established.
\editorialnote{A gauge-invariant Majorana mass for the active Standard Model
neutrino is not a renormalizable bare term before electroweak breaking. At
low energy it can arise from the dimension-five Weinberg operator
\((LH)(LH)/\Lambda\), or through a seesaw involving heavy sterile neutrinos.
Whether neutrinos are Dirac or Majorana particles remains experimentally
unsettled.}

\section{Why the Familiar Spectrum Clusters Near the Electroweak Scale}

\subsection{The common scale and its important qualifications}

\begin{classroomqa}
\audiencequestion{Why was the Higgs mass expected to be reachable at the
weak scale?\lecturetimestamp{01:14:48}}
\lecturerresponse{The curvature of the Higgs potential sets the radial
Higgs mass, while the same vacuum scale \(f\) sets the \(W\) and \(Z\)
masses. Without an extreme scalar self-coupling, one therefore expected the
Higgs mass to be of the same broad order as the weak scale. The observed
\(125\ \mathrm{GeV}\) value indeed lies there.}
\end{classroomqa}

The lecture's striking summary is that the elementary masses of the minimal
Standard Model trace back to one broken vacuum:
\begin{align}
  m_W&=\frac12gv,\\
  m_Z&=\frac12\sqrt{g^2+g'^2}\,v,\\
  m_f&=\frac1{\sqrt2}y_fv.
  \label{eq:l09-electroweak-mass-summary}
\end{align}
If \(v\) vanished while the dimensionless couplings were held fixed, the
weak gauge bosons and charged elementary fermions would be massless.

That statement must not be enlarged beyond its domain. The photon and
gluons remain massless because an unbroken gauge symmetry protects them.
Most of the proton and neutron mass is QCD binding energy, not the sum of
their quark Higgs masses. Neutrino mass requires an extension of the
minimal model. The Higgs mass itself depends on the scalar potential and is
not simply another Yukawa mass.

\subsection{Particles whose mass need not wait for symmetry breaking}

Suppose a new fermion had left- and right-chiral components with identical
gauge quantum numbers. Then
\(\bar\Psi_L\Psi_R+\mathrm{h.c.}\) would already be gauge invariant. Such a
\emph{vectorlike} fermion could possess a bare mass \(M\) unrelated to
\(v\):
\begin{equation}
  \mathcal L_{\mathrm{vectorlike}}
  =-M\bar\Psi\Psi.
  \label{eq:l09-vectorlike-mass}
\end{equation}
There is no electroweak symmetry obstruction to \(M\gg v\).

This motivates the lecture's broad naturalness argument. Particles protected
by gauge or chiral symmetry remain massless until the Higgs vacuum breaks
that protection, so their masses are tied to the relatively small scale
\(v\). Particles with no such protection may naturally be much heavier and
therefore absent from the observed low-energy spectrum. What we see may be
a selection effect: the accessible particles are the ones whose masses were
forced to wait for electroweak symmetry breaking.

\begin{classroomqa}
\audiencequestion{Would particles with independent, much larger masses play
no role in reality?\lecturetimestamp{01:19:39}}
\lecturerresponse{They would play little role in ordinary low-energy
laboratory physics because they are hard to produce, but they could still
shape high-energy history or cosmology. The lecture offers dark matter as a
possible place where a heavier hidden spectrum might first reveal itself.}
\end{classroomqa}

The 2010 discussion reflects the then-popular expectation of weak-scale or
somewhat heavier dark matter. That remains a useful model class, not an
established fact. Dark matter has been detected gravitationally, but its
particle identity and the origin of its mass remain unknown.

\section{Closing Questions: Cosmic Rays, Inertia, Gravity, and Chirality}

\subsection{Mass is not high laboratory energy}

\begin{classroomqa}
\audiencequestion{What is the mass of a cosmic ray?
\lecturetimestamp{01:20:57}}
\lecturerresponse{A cosmic ray may be a proton, an atomic nucleus, an
electron, or a photon; its invariant mass is the mass of that particle or
system. What is spectacular is its laboratory energy, which for the highest
events reaches roughly \(10^{20}\) to \(10^{21}\ \mathrm{eV}\). High energy
is not the same as high invariant mass.}
\end{classroomqa}

A cosmic ray strikes a target particle that is approximately at rest. For a
projectile proton of laboratory energy \(E\gg m_p\),
\begin{equation}
  s=(p_1+p_2)^2
  \simeq2m_pE+2m_p^2,
  \qquad
  \sqrt{s}\simeq\sqrt{2m_pE}.
  \label{eq:l09-fixed-target-s}
\end{equation}
Only the square root of the projectile energy becomes center-of-mass energy.
For equal head-on beams of energy \(E_{\rm beam}\),
\begin{equation}
  \sqrt{s}\simeq2E_{\rm beam}.
  \label{eq:l09-collider-s}
\end{equation}
This is why a collider is far more efficient than a fixed target at
converting controlled beam energy into new-particle mass.

\begin{classroomqa}
\audiencequestion{Why can CERN's collider be a more potent probe than a
single much more energetic particle hitting matter at rest?
\lecturetimestamp{01:21:50}}
\lecturerresponse{In a fixed-target collision, momentum conservation carries
most of the laboratory energy forward in the motion of the final state. In
a symmetric collider the total spatial momentum can vanish, leaving the
beam energy available as center-of-mass energy. The very highest cosmic rays
can still exceed the LHC in \(\sqrt{s}\), but they arrive unpredictably and
cannot provide collider luminosity or a controlled detector environment.}
\end{classroomqa}

\subsection{Rest mass and resistance to acceleration}

\begin{classroomqa}
\audiencequestion{Are inertia and mass essentially the same thing?
\lecturetimestamp{01:22:41}}
\lecturerresponse{At rest, invariant mass is both rest energy through
\(E_0=mc^2\) and the parameter governing the low-velocity response to a
force. Special relativity explains why those definitions belong to one
quantity. For moving bodies it is clearer to keep invariant mass fixed and
describe acceleration using relativistic energy and momentum.}
\end{classroomqa}

The invariant relation
\begin{equation}
  E^2=p^2c^2+m^2c^4
  \label{eq:l09-energy-momentum}
\end{equation}
contains both meanings. Near rest,
\begin{equation}
  E=mc^2+\frac{p^2}{2m}+\cdots,
\end{equation}
so the same \(m\) that gives the rest energy appears in the kinetic response.
At high speed the relation between force and acceleration depends on whether
the force is parallel or perpendicular to the velocity. This directional
dependence is why modern treatments avoid a velocity-dependent
\emph{relativistic mass} and reserve mass for the Lorentz invariant \(m\).

\subsection{What gravity couples to}

\begin{classroomqa}
\audiencequestion{Once the Higgs gives a particle mass, does it then begin
to couple to gravity?\lecturetimestamp{01:23:58}}
\lecturerresponse{Its rest energy contributes to gravity, but mass is not
the whole source. Photons are massless and are nevertheless deflected by
gravity. In general relativity the source is the stress--energy tensor:
energy density, momentum flow, pressure, and stress. For slowly moving
matter the dominant contribution is its rest-mass energy.}
\end{classroomqa}

The bending of light near the Sun makes the point operationally. A photon
responds to spacetime curvature despite having \(m=0\), and its energy and
momentum also contribute to the gravitational field. For two nearly static
massive bodies, energy is dominated by \(mc^2\), so Newtonian gravity can be
written in terms of their masses without displaying the more general
stress--energy statement.

\begin{classroomqa}
\audiencequestion{Is a high-energy particle harder to accelerate in the same
way in every direction?\lecturetimestamp{01:25:20}}
\lecturerresponse{No. The change of momentum for a force parallel to the
motion differs from the change for a transverse force. It is correct that a
more energetic particle is harder to deflect, but assigning it one enlarged
moving mass hides this directional structure.}
\end{classroomqa}

\subsection{An electron at rest and the promised distinction}

\begin{classroomqa}
\audiencequestion{How can an electron at rest be called left- or
right-handed?\lecturetimestamp{01:26:32}}
\lecturerresponse{It cannot be assigned helicity at rest because helicity
uses the direction of momentum. A rest spinor can still be decomposed into
left- and right-chiral components. For a massive electron both components
are present, and the mass term couples them.}
\end{classroomqa}

\begin{classroomqa}
\audiencequestion{If an electron at rest has no helicity, can it still
participate in a weak interaction?\lecturetimestamp{01:27:18}}
\lecturerresponse{Yes. The weak interaction couples to chirality, not to a
frame-dependent helicity label. At high energy helicity is a convenient
approximation to chirality; near rest that shortcut fails. This is the
subtle distinction the lecture had postponed.}
\end{classroomqa}

\section{Historical End-of-Lecture Questions}

\subsection{The pre-discovery Higgs window}

In 2010 the Higgs had not yet been seen. Precision electroweak observables,
especially loop-sensitive relations among the \(W\), \(Z\), top, and other
parameters, favored a comparatively light Higgs within the Standard Model.
The lecture quotes a likely upper region near \(160\ \mathrm{GeV}\), well
below \(v\) but not separated from it by an order of magnitude.

\begin{classroomqa}
\audiencequestion{What evidence then suggested that the Higgs was lighter
than the \(250\ \mathrm{GeV}\) vacuum scale?\lecturetimestamp{01:28:28}}
\lecturerresponse{High-precision measurements could be compared with
radiative corrections whose values depend logarithmically on the Higgs
mass. The resulting global fits favored a light Higgs. This was indirect
evidence, not a direct mass measurement; the later direct value is about
\(125\ \mathrm{GeV}\).}
\end{classroomqa}

\begin{classroomqa}
\audiencequestion{Could several minima of the Higgs potential produce
different values of \(f\), perhaps tied to different particle generations?
\lecturetimestamp{01:29:28}}
\lecturerresponse{Different minima can support very different low-energy
physics, but there is no simple inference from the observed generations to
multiple Higgs minima. The lecture defers the broader subject of multiple
vacua to a later course rather than pretending to settle it here.}
\end{classroomqa}

\begin{classroomqa}
\audiencequestion{Was the Higgs expected to be the first new particle found
at the LHC?\lecturetimestamp{01:31:32}}
\lecturerresponse{Not necessarily. Its visibility depends on production
rates and decay channels, some of which can be difficult or invisible.
Before LHC data, supersymmetric or other new particles could plausibly have
appeared first. Historically, the Higgs was discovered and no supersymmetric
particle has yet been established.}
\end{classroomqa}

\subsection{When a neutral particle is its own antiparticle}

\begin{classroomqa}
\audiencequestion{Do the photon and gluon lack antiparticles because they
are massless?\lecturetimestamp{01:32:12}}
\lecturerresponse{No. Masslessness is not the criterion. A photon is its own
antiparticle because no conserved charge distinguishes a photon from an
antiphoton. Each gluon is also a self-conjugate gauge quantum in the
appropriate color representation, although gluons themselves carry color
and the color labels mix under charge conjugation. A neutron is electrically
neutral but has a distinct antineutron because baryon number and its quark
content distinguish them.}
\end{classroomqa}

The neutrino returns as the subtle case. It has weak interactions, and the
ultrarelativistic states called neutrino and antineutrino interact with
opposite helicities. If an exact lepton number distinguishes them, they can
form a Dirac particle with a distinct antiparticle. If lepton number is not
exact and the mass is Majorana, the massive particle is its own
antiparticle even though its two helicity states behave differently in weak
processes.

\begin{classroomqa}
\audiencequestion{Do weak charge and lepton number prove that neutrinos and
antineutrinos must be distinct?\lecturetimestamp{01:32:46}}
\lecturerresponse{Not conclusively. Weak interactions distinguish the
helicity states, but a Majorana mass can combine a neutrino field with its
charge conjugate. The decisive issue is whether an exact conserved quantum
number forbids that mixing. Total lepton number is not imposed as an exact
gauge symmetry by the Standard Model.}
\end{classroomqa}

Ordinary flavor oscillations require a separate caution.
\editorialnote{Ordinary observed flavor oscillations mix
\(\nu_e,\nu_\mu,\nu_\tau\) of the same neutrino or antineutrino sector; they
do not by themselves demonstrate neutrino--antineutrino conversion.
Lepton-number-violating processes such as neutrinoless double-beta decay
would provide direct evidence for Majorana mass, but none has yet been
confirmed.}

\begin{classroomqa}
\audiencequestion{Does the Standard Model explain why electric charges are
quantized?\lecturetimestamp{01:33:39}}
\lecturerresponse{The lecture acknowledges the question but does not develop
it. Within the Standard Model, gauge invariance and anomaly cancellation
strongly constrain consistent hypercharge assignments; a deeper explanation
often invokes grand unification or other ultraviolet structure.}
\end{classroomqa}

\section{What the Lecture Has Established}

The gauge-boson and fermion stories now line up.

\begin{enumerate}
  \item A derivative-only gauge theory has no energy gap for a homogeneous
  vector potential. Gauge symmetry forbids a bare vector mass.
  \item A charged scalar with nonzero vacuum magnitude contributes
  \(g^2f^2A_\mu A^\mu\). Its angular Goldstone mode becomes the longitudinal
  polarization of a massive weak gauge boson.
  \item A Dirac mass couples \(\psi_L\) and \(\psi_R\). Because the weak
  interaction assigns them different gauge quantum numbers, a bare
  left--right mass term is forbidden.
  \item A Yukawa interaction inserts the Higgs field between the two
  chiralities and is gauge invariant. Replacing the Higgs by its vacuum
  value gives \(m_f=y_ff\), while retaining its fluctuation gives the
  physical Higgs--fermion coupling.
  \item Neutrino mass obeys the same left--right logic but admits both Dirac
  and Majorana realizations because the neutrino is electrically neutral.
  \item Particles without gauge or chiral protection could possess masses
  unrelated to the electroweak scale and may lie far above present
  experiments.
\end{enumerate}

The unifying statement is therefore not that the Higgs is a substance that
mechanically drags on particles. It is that the electroweak vacuum has a
nonzero scalar order parameter. That order parameter changes which
gauge-invariant terms are possible in the low-energy equations. For vectors
it opens a zero-momentum gap; for fermions it permits left--right mixing.
The different masses then measure different couplings to the same broken
vacuum.
